\begin{abstract}
Scientific workflows rely on schedulers to reserve memory before tasks
start, balancing execution failures against excess reservations. Consumption-Aware Memory
Prediction (\toolname), introduced in our prior WORKS paper, augments task
metadata and memory history with measured and predicted data consumption.
This study examines the value and cost of that feedback through component
analyses, selective-auditing replay, and new Nextflow/Slurm executions.
The component results expose tradeoffs between prediction error, allocation
coverage, and auditing cost rather than a uniformly dominant configuration.
In a live Bowtie2 evaluation with 6,000 tasks per stage, an additional pass
over the original inputs reduces aggregate CAMP requests by 75.34\%.
On a different ENA study, initial-request coverage falls from 99.58\% on
repeated inputs to 94.15\%, despite continued online learning.
Branching directly from the same cold checkpoint, without the extra pass,
increases requests by 7.62\% and coverage to 95.78\%.
All tasks complete with no recorded OOMs under a separate 128~MiB per-task
execution allowance; these observations do not establish the sufficiency
of CAMP requests alone. The results support evaluating consumption feedback
through its coverage--reservation tradeoff, while distinguishing online
adaptation from frozen-model generalization and allocation estimates from
end-to-end resource savings.
\end{abstract}
\begin{IEEEkeywords}
scientific workflows, online memory allocation, consumption feedback,
eBPF, selective auditing
\end{IEEEkeywords}
